The general projectivity problem admits a presentation-independent support
condition. For a left $Q$-semimodule $L$, let $\mathcal A_L$ be the
preadditive action category of Section~\ref{sec:ringoid}: its objects are
$l\in L$, with transition arrows $t_{l,k}$ and scalar arrows
$a_{q,x,l}:l\to ql$, subject to the module relations displayed there.
Call a $Q$-linear map $\phi:L\to Q$ \emph{admissible at $l$} if
$\phi(k)l\leq k$ for every $k\in L$. Define
\[
 \mathcal J_L(l,k)=
 \left\langle
 t_{\phi(k)l,k}a_{\phi(k),x,l}:
 \phi\text{ admissible at }l,\ x\in F(\phi(k))
 \right\rangle_{\Z}.
\]
\begin{theorem}[canonical projectivity ideal]\label{thm:canonical-ideal}
For every unital inverse semiring $T=(Q,F)$, these groups form a
two-sided ideal of $\mathcal A_L$. A $T$-semimodule with support $L$ and
amplitude $G$ is finitely generated projective if and only if
\begin{gather*}
 L\text{ is finitely generated projective over }Q,\qquad
 G\text{ is finitely generated projective over }\mathcal A_L,\\
 \mathcal J_LG=G.
\end{gather*}
Equivalently, the last condition is
$(\mathcal A_L/\mathcal J_L)\otimes_{\mathcal A_L}G=0$.
For a fixed identified support $L$, its projectives are thus exactly the
finite projective $\mathcal A_L$-modules generated by this ideal.
Support isomorphisms transport the category, ideal and amplitude data.
No commutativity, finiteness or injectivity assumption on $(Q,F)$ is required.
\end{theorem}
The ideal depends only on the scalar diagram and the support. It replaces
the existential support-section test of Theorem~\ref{thm:ringoid-projective}
by ordinary module projectivity and generation by one specified ideal.
Classifying arbitrary projective $Q$-supports and projectives over the
resulting ordinary action categories remains a separate problem.
